% =====================================================================
% HOJA 9 (ix) – ABSTRACT AND KEYWORDS
% =====================================================================
\begin{center}
    \vspace*{0.5cm}
    {\bfseries ABSTRACT}
\end{center}
\vspace{0.8cm}

\justifying
\selectlanguage{english}
\noindent
The present project aims to develop ``Kiddy Quiz,'' an interactive web application designed for the evaluation and reinforcement of mathematics for elementary school students. For the theoretical foundation and requirements gathering, a field research study was conducted, which included interviews with 5 education professionals and 2 psychopedagogues. This allowed for the identification of cognitive and anxiety barriers present in traditional evaluation processes. The software development was managed under the Extreme Programming (XP) agile methodology, ensuring code quality and robustness. The system architecture was built using Angular for frontend development, NestJS for backend logic, and PostgreSQL as the relational database engine. As a primary innovation, the platform integrates pictogram-based interfaces to facilitate autonomous navigation for children, as well as the implementation of Artificial Intelligence (Gemini) to provide motivational feedback after each evaluation. System validation was carried out using the System Usability Scale (SUS) and the Technology Acceptance Model (TAM), demonstrating an ``Excellent'' usability category and an almost unanimous intention of use by the students. Furthermore, to guarantee data protection and proper role-based access control, security measures were implemented using JWT token authentication. As a result, an inclusive technological tool has been consolidated, optimizing teachers' academic management, reducing student frustration, and representing a significant contribution to the modernization of childhood educational assessment.

\vspace{0.5cm}

\noindent\textbf{Keywords:} web application, elementary education, mathematics, pictograms, artificial intelligence, usability, Kiddy Quiz.
\selectlanguage{spanish}

\newpage
